\documentclass[aspectratio=169,usenames,dvipsnames,compress]{beamer}

\usepackage{scrextend}
\usepackage[T1]{fontenc}
\usepackage{lmodern}
\usepackage{comment}
\usepackage{amsmath}
\usepackage{caption}
\usepackage{booktabs}
\usepackage{array}

\usepackage{enumerate} 
\usepackage{tikz}
\usetikzlibrary{arrows.meta, positioning, shapes.geometric}


% ********************************************************

\makeatletter
\usepackage[authoryear]{natbib}
%\renewcommand{\bibsection}{}
\setbeamertemplate{frametitle continuation}{}

\usepackage{enumitem}
\usepackage{array}
\usepackage{tikz}
\def\checkmark{\tikz\fill[scale=0.4](0,.35) -- (.25,0) -- (1,.7) -- (.25,.15) -- cycle;} 
\definecolor{myblue}{RGB}{0, 0, 255}
\definecolor{lightyellow}{RGB}{255,255,224} % Define a cor lightyellow
\definecolor{lightred}{RGB}{255,182,193} % Define a cor lightred
\usetikzlibrary{positioning}

\mode<presentation>
    {
    \usetheme{Warsaw} 
    %\usetheme{Madrid}
    \usefonttheme{professionalfonts} 
    \setbeamertemplate{itemize item}{$\blacksquare$}
    \setbeamertemplate{itemize subitem}{$\blacktriangleright$}
    }

% Personalização do Warsaw Theme

% -------------------------
% Beamer Fontes
% -------------------------
\changefontsizes{10pt}
\setbeamerfont{title}{size=\small}
\setbeamerfont{author}{size=\small}
\setbeamerfont{frametitle}{size=\small}
\setbeamerfont{block title}{size=\small}
\setbeamerfont{footline}{series=\tiny}

% How to insert page number in Beamer navigation symbols?

\addtobeamertemplate{navigation symbols}{}{%
    \usebeamerfont{footline}%
    \usebeamercolor[fg]{footline}%
    \hspace{1em}%
    \insertframenumber/\inserttotalframenumber
}

\setbeamercolor{section in head/foot}{fg=white}
\setbeamercolor{subsection in head/foot}{fg=white}
%\setbeamercolor{footline}{fg=white, bg=darkgray}
% Colocar indices de avaço no head da pagina

\defbeamertemplate*{headline}{}
{
    \begin{beamercolorbox}[colsep=1.5pt]{upper separation line head}
    \end{beamercolorbox}
    \begin{beamercolorbox}{section in head/foot}
    \vskip2pt\insertnavigation{\paperwidth}\vskip4pt
    \end{beamercolorbox}
    \begin{beamercolorbox}[colsep=1.5pt]{middle separation line head}
    \end{beamercolorbox}
    \begin{beamercolorbox}[colsep=1.5pt]{lower separation line head}
    \end{beamercolorbox}
}

%\usepackage{beamerhighlight}
%\usepackage{pgf, pgfpages}
%\usepackage{colortbl}
%\usepackage{color}
%\usepackage{pdfpages}
%\usepackage[absolute,overlay]{textpos}
%\usepackage{url}
%\usepackage{graphicx}

% Acentos em modo matemático
\usepackage{amstext}
\usepackage{amsmath}
\usepackage{enumitem}
%\usepackage{wasysym}
%usepackage[squaren]{SIunits}


% Numerar as figuras em Caption
\setbeamertemplate{caption}[numbered]

 %Justificar Texto
\usepackage{ragged2e}
\usepackage{etoolbox}
\apptocmd{\frame}{}{\justifying}{} % Allow optional arguments after frame.

%\definecolor{Azul}{RGB}{0,0,250}
\definecolor{Azul}{rgb}{0,0,1}
\definecolor{Verde}{rgb}{0,1,0}
\definecolor{Red}{rgb}{1,0,0}
\definecolor{Orange}{rgb}{0.9,0.4,0.8}
\definecolor{aqua}{rgb}{0.0, 1.0, 1.0}
\definecolor{aquamarine}{rgb}{0.5, 1.0, 0.83}
\definecolor{arsenic}{rgb}{0.23, 0.27, 0.29}

\usepackage{xcolor}

%\usepackage[dvipsnames]{xcolor}
\usepackage{multirow}

% Numeração arabiga em Footnote Beamer
\renewcommand\thempfootnote{\arabic{mpfootnote}}
% Fonte do footnote
\setbeamerfont{footnote}{size=\tiny}


\usepackage[hang]{footmisc}
\setlength{\footnotemargin}{3mm}

%\def\setnotebodyfont
%{\switchtobodyfont[tiny]\setupinterlinespace[15]} 


% TColorBox
\usecolortheme{whale}
\usecolortheme{orchid}
%\usepackage{tcolorbox}
\usepackage[most]{tcolorbox}
%\tcbset{colback=red!5!white,colframe=red!75!black}
\usepackage{etoolbox}
\AtBeginEnvironment{tcolorbox}{\footnotesize}

\usepackage{multicol,lipsum}

% Desabilitra avisos
\usepackage{silence}
%Disable all warnings issued by latex starting with "You have..."
\WarningFilter{latex}{You have requested}


% Alterar simbolos no Itemize

%\renewcommand{\labelitemi}{$\bullet$}
%\renewcommand{\litemii}{$\cdot$}
%\renewcommand{\litemiii}{$\diamond$}
%\renewcommand{\litemiv}{$\ast$}

% Configuraçao do Itemize
\def\labelitemi{$\bullet$}
\usepackage{enumitem}
\setlist[itemize]{align=parleft,left=-2pt..1em}



\usepackage{amssymb}% \varnothing
\usepackage{latexsym}

\usepackage[geometry]{ifsym}
\usepackage{amsmath}
\usepackage{amssymb}
\usepackage{latexsym}
\usepackage{relsize}
\usepackage{color,graphicx}
\usepackage{xfrac}


%\usepackage{boisik}
%\usepackage{lmodern}% http://ctan.org/pkg/lm
%\usepackage{amsmath}

\usepackage{stmaryrd}
%\SetSymbolFont{stmry}{bold}{U}{stmry}{m}{n}

\newtheorem{mydescription}{Descrição}[section]
\usepackage{calc}
\usepackage{enumitem}
\SetLabelAlign{parright}{\parbox[t]{\labelwidth}{\raggedleft#1}}

\usepackage{listings}
\usepackage{array,tabularx}
\usepackage{colortbl}
%\usepackage[table]{xcolor}

% FIGURE

\usepackage[font=scriptsize,skip=2pt]{caption}

\usepackage{tabularx} 
\usepackage{xcolor}
\newcommand{\cor}{\textcolor}
\usepackage{media9}
\usepackage{comment}

%LaTeX Font Warning: Font shape `OT1/cmr/m/n
%Timeline for How to remove the warnings “Font shape `
\usepackage{anyfontsize}
\usepackage{caption}
\usepackage{subcaption}
\usepackage{framed}

%\usepackage[lighttt]{lmodern}
\usepackage{listings}
\lstset{basicstyle=\ttfamily, keywordstyle=\bfseries}

\usepackage{multicol}
\usepackage{parskip}
\usepackage{pxfonts}
\usepackage{graphics}
\usepackage{pstricks}

% Texto justificado
\usepackage{ragged2e}
\usepackage{colortbl}
\usepackage{comment}

% ---------------------------------------------------
% Set font size for footnotes

%\makeatletter
%\renewcommand\footnotesize{%
%   \@setfontsize\scriptsize\@ixpt{7}%
%   \abovedisplayskip 8\p@ \@plus2\p@ \@minus4\p@
%   \abovedisplayshortskip \z@ \@plus\p@
%   \belowdisplayshortskip 4\p@ \@plus2\p@ \@minus2\p@
%   \def\@listi{\leftmargin\leftmargini
%               \topsep 4\p@ \@plus2\p@ \@minus2\p@
%               \parsep 2\p@ \@plus\p@ \@minus\p@
%               \itemsep \parsep}
%   \belowdisplayskip \abovedisplayskip
%}
%\makeatother

% Set font size for footnotes
%\renewcommand{\footnotesize}{\scriptsize}

%You can use arbitrary font sizes within LaTeX, not only the default ones 
% (\normalsize, \Large, \footnote size, &c.). The command for a font size of 9pt would be:

\renewcommand{\footnotesize}{\fontsize{6pt}{9pt}\selectfont}

%The second value is value for the interline skip for that size. 
%Usually it is about 20 % more than the nominal size.
% ---------------------------------------------------
%\AtBeginDocument{\fontsize{8}{10}\selectfont}

\UseRawInputEncoding
%\usepackage{ae,lmodern}
% Tabelas TcolorBox


\makeatother

%\usepackage{tcolorbox}
\tcbuselibrary{skins}
%\tcbset{beamer without segmentation/.style={%
%  beamer,segmentation code=,  interior titled code={{\tcb@spec{beamer@color}\tcb@drawwithtitle@path}\tcb@drawspec@T},  interior code={{\tcb@spec{beamer@color}\tcb@drawwithouttitle@path}}}}


\newcommand{\PR}[1]{\ensuremath{\left[#1\right]}}
\newcommand{\PC}[1]{\ensuremath{\left(#1\right)}}
\newcommand{\CH}[1]{\ensuremath{\left\{#1\right\}}}


%\newcommand\scaledfrac[2]{\scalebox{2}{\sfrac{#1}{#2}}}

%\newcommand{cmd}[args][opt]{def}
%\newtheorem{Exemplo}{Exemplo}[section]
% Closed” (square) root symbol

\usepackage{letltxmacro}
\makeatletter
\let\oldr@@t\r@@t
\def\r@@t#1#2{%
\setbox0=\hbox{$\oldr@@t#1{#2\,}$}\dimen0=\ht0
\advance\dimen0-0.2\ht0
\setbox2=\hbox{\vrule height\ht0 depth -\dimen0}%
{\box0\lower0.4pt\box2}}
\LetLtxMacro{\oldsqrt}{\sqrt}
\renewcommand*{\sqrt}[2][\ ]{\oldsqrt[#1]{#2}}
\makeatother



\usepackage[brazilian]{babel}
%\usepackage{lmodern}			% Usa a fonte Latin Modern
\usepackage[utf8]{inputenc}
\usepackage[T1]{fontenc}    % Selecao de codigos de fonte.



%\definecolor{Azul}{RGB}{0,0,250}
\definecolor{Blue}{rgb}{0,0,1}
\definecolor{Green}{rgb}{0,1,0}
\definecolor{Red}{rgb}{1,0,0}
\definecolor{Magenta}{rgb}{1,0,0.56}
\definecolor{Cyan}{rgb}{0,1,1}
\definecolor{Yellow}{rgb}{1,1,0}
\definecolor{White}{rgb}{1,1,1}

\lstdefinelanguage{Scala}
{
 basicstyle=\ttfamily\footnotesize
}


\makeatletter
\newcommand\titlegraphicii[1]{\def\inserttitlegraphicii{#1}}
\titlegraphicii{}
\setbeamertemplate{title page}
{
  \vbox{}
   {\usebeamercolor[fg]{titlegraphic}\inserttitlegraphic\hfill\inserttitlegraphicii\par}
  \begin{centering}
    \begin{beamercolorbox}[sep=8pt,center]{institute}
      \usebeamerfont{institute}\insertinstitute
    \end{beamercolorbox}
    \begin{beamercolorbox}[sep=8pt,center]{title}
      \usebeamerfont{title}\inserttitle\par%
      \ifx\insertsubtitle\@empty%
      \else%
        \vskip0.25em%
        {\usebeamerfont{subtitle}\usebeamercolor[fg]{subtitle}\insertsubtitle\par}%
      \fi%     
    \end{beamercolorbox}%
    \vskip1em\par

 %\begin{beamercolorbox}[sep=8pt,center]{date}
  %\begin{beamercolorbox}[sep=1pt,center]{}
      %\usebeamerfont{date}\insertdate
    %\end{beamercolorbox}%\vskip0.5em
    \begin{beamercolorbox}[sep=2pt,center]{author}
      \usebeamerfont{author}\insertauthor
    \end{beamercolorbox}
  \end{centering}
  %\vfill
}


\newcommand{\code}{\ttfamily}
\setcounter{page}{1}
\usepackage{graphics} 
\usepackage{ifpdf}
\ifpdf
  \DeclareGraphicsRule{*}{mps}{*}{}
\fi


\newcommand{\Tab}{\hspace{0.5cm}}
\newcommand{\TAB}{\hspace{1cm}}


%\input{Tools/C++.tex}

\usepackage{listings,newtxtt}
\usepackage{xcolor}

\definecolor{commentsColor}{rgb}{0.497495, 0.497587, 0.497464}
\definecolor{keywordsColor}{rgb}{0.000000, 0.000000, 0.635294}
\definecolor{stringColor}{rgb}{0.558215, 0.000000, 0.135316}
\definecolor{arylideyellow}{rgb}{0.91, 0.84, 0.42}
\definecolor{banana}{rgb}{0.98, 0.91, 0.71}


\lstset{ %
  language=Python,                 % the language of the code (can be overrided per snippet)
  basicstyle=\ttfamily,            % fonts that are used for the code
  backgroundcolor=\color{green!10},% choose the background color; 
  breakatwhitespace=false,         % sets if automatic breaks should only happen at whitespace
  breaklines=true,                 % sets automatic line breaking
  captionpos=t,                    % sets the caption-position to bottom
  commentstyle=\color{commentsColor}\textit,    % comment style
  deletekeywords={...},            % if you want to delete keywords from the given language
  extendedchars=true,              % lets you use non-ASCII characters; does not work with UTF-8
  escapeinside={(*}{*)},           % LaTeX
  frame=t,                         % adds a frame around the code
  keepspaces=true,                 % keeps spaces in text
  keywordstyle=\color{keywordsColor}\bfseries,       % keyword style
  otherkeywords={*,...},           % if you want to add more keywords to the set
  numbers=right,                   % where to put the line-numbers; possible values are (none, left, right)
  numbersep=5pt,                   % how far the line-numbers are from the code
  %xleftmargin=.25in,               % indentação
  %xrightmargin=.25in}
  numberstyle=\tiny\color{commentsColor}, % the style that is used for the line-numbers
  rulecolor=\color{black},         % if not set, the frame-color may be changed on line-breaks within not-black text (e.g. comments (green here))
  showspaces=false,                % show spaces everywhere adding particular underscores; it overrides 'showstringspaces'
  showstringspaces=false,          % underline spaces within strings only
  showtabs=false,                  % show tabs within strings adding particular underscores
  stepnumber=1,     % the step between two line-numbers. If it's 1, each line will be numbered
  stringstyle=\color{stringColor}, % string literal style
  tabsize=1,                       % sets default tabsize to 2 spaces
  title=\lstname,                  % show the filename of files included with \lstinputlisting; also try caption instead of title
  columns=fixed                    % Using fixed column width (for e.g. nice alignment)
  %extendedchars=\true,
  %inputencoding=utf8x,
  }


\lstdefinestyle{interfaces}
  {
  float=tp,
  floatplacement=tbp,
  abovecaptionskip=-5pt
  }

\lstdefinelanguage{Python}
  {
  keywords={print, import, True, False, plot, Mesh, array, plt, def, return, type, lambda, map, list, while, if , elif, else, for, in, __name__ },
  keywordstyle=\color{blue}\bfseries,
  ndkeywords={boolean, throw, implements, import, this, not, int, float, while, open,close, read, readline, redalines,write,writelines, seek, split, len},
  ndkeywordstyle=\color{red}\bfseries,
  identifierstyle=\color{black},
  sensitive=false,
  comment=[l]{//},
  morecomment=[s]{/*}{*/},
  commentstyle=\color{purple}\ttfamily,
  stringstyle=\color{red}\ttfamily,
  morestring=[b]',
  morestring=[b]"
}



\lstset{
language=Python,   % R code
literate=	{á}{{\'a}}1
			{à}{{\`a}}1
			{ã}{{\~a}}1
{é}{{\'e}}1
{ê}{{\^e}}1
{í}{{\'i}}1
{ó}{{\'o}}1
{õ}{{\~o}}1
{ú}{{\'u}}1
{ü}{{\"u}}1
{ç}{{\c{c}}}1
}

\renewcommand{\lstlistingname}{\scriptsize Código}

\usepackage{inconsolata}

%\input{Tools/Haskell.tex}
%\input{Tools/Hugs.tex}
\usepackage{microtype}
\usepackage{roboto}
\usepackage{tcolorbox}
\tcbuselibrary{skins}
\usepackage{tcolorbox}


\newtcolorbox[auto counter, number within=section]{Constative}[3]{%
  enhanced,
  size=minimal, left=3.0cm, top=3mm, toprule=2pt,
  colframe=#3, colbacktitle=#3, colback=white,
  fontupper=\sffamily\scriptsize, title=#2,
  attach boxed title to top left={yshift=-\tcboxedtitleheight},
  center title, minipage boxed title=2.5cm,
  fonttitle=\sffamily\robotocondensed\bfseries\small,
  boxed title style={%
    size=minimal, boxsep=4pt, bottom=5mm,
    frame code={%
      \path[fill=tcbcolback]
      (frame.north west) -- ([yshift=2mm]frame.south west) --
      ([xshift=2mm]frame.south west) -- (frame.south east) --
      (frame.north east) -- cycle;
      \path[fill=white, draw=tcbcolback, line width=1pt]
      ([xshift=-15mm-0.5pt]frame.south east) circle [radius=5mm]
      node [font=\sffamily\small] {#1};
    },
    interior engine=empty
  },
  frame code={%
    \path[fill=tcbcolframe](frame.north west) rectangle ([yshift=-3pt]frame.north east);
  }
}

% ****** Intertext like command in enumerate environment?

\usepackage{enumitem}
\usepackage{exercises}

% ---------------------------------
% Floor Ceil Math
% ---------------------------------
\usepackage{relsize}

% -------------------------
% Comandos matematicos
% ***************************
%\newcommand{\BigMath}[1] {\mathlarger{\mathlarger{\mathlarger{#1}}}} 
%\newcommand{\LargeMath}[1] {\mathlarger{\mathlarger{#1}}} 
%\newcommand{\largeMath}[1] {\mathlarger{#1}}
\newcommand{\BigMath}[1] {\mathlarger{\mathlarger{\mathlarger{\mathlarger{\mathlarger{\mathlarger{#1}}}}}}} 
\newcommand{\LargeMath}[1] {\mathlarger{\mathlarger{#1}}} 
\newcommand{\largeMath}[1] {\mathlarger{#1}}

% Símbolos matemáticos

\newcommand{\SMath}[2]{\scalebox{#1}{#2}}
%\newcommand{\largeMath}[1] {\mathlarger{#1}}
%$\scalebox{1.2}{f(3) = 32 + 7 = 16}$ \\

\usepackage{enumerate}
\usepackage{mathtools}
\DeclarePairedDelimiter\ceil{\lceil}{\rceil}
\DeclarePairedDelimiter\floor{\lfloor}{\rfloor}

% ---------------------------------
% Fancy Atividades
% ---------------------------------

\usepackage{microtype}
\usepackage{roboto}
\usepackage{tcolorbox}
\tcbuselibrary{skins}

\usepackage{tikz}
%\def\checkmark {\tikz\fill[scale=0.4](0,.35) -- (.25,0) -- (1,.7) -- (.25,.15) -- cycle;} 
\def\checkmark{\tikz\fill[scale=0.4](0,.35) -- (0.25,0) -- (0.5,0.7) -- (0.25,0.15) -- cycle;} 

%Package microtype Warning: Unable to apply patch
\let\CheckCommand\providecommand
\usepackage{microtype}

% Colored itemize
%\usepackage{beamerhighlight}

% description
\usepackage{calc}  
\usepackage{enumitem}

\newcommand*{\listnum}{\sffamily\robotoblack\color{cyan!10! black}}
\newcommand*{\textlistnum}[1]{\begingroup\listnum #1\endgroup}

\setlist[enumerate,1]{label=\listnum\arabic*}
\setlist[enumerate,2]{label=\listnum\alph*, leftmargin=1.5em}
\newcommand\setItemnumber[1]{\setcounter{enum\romannumeral\@enumdepth}{\numexpr#1-1\relax}}
\SetLabelAlign{parright}{\parbox[t]{\labelwidth}{\raggedleft#1}}

% ---------------------------------
% Fancy Atividades
% ---------------------------------

\newcommand{\BR}{\rowcolor{blue}}
\newcommand{\GR}{\rowcolor{green!85!yellow}}
\newcommand{\Aula}{\rowcolor{green}}
\newcommand{\Feriado}{\rowcolor{blue!45!white}}
\newcommand{\Prova}{\rowcolor{red}}
\newcommand{\Anuncios}{\rowcolor{gray!60!black}}
\newcommand{\CR}{\rowcolor{black!20!white}}
\newcommand{\Projeto}{\rowcolor{yellow!12!green}}
\newcommand{\WT}[1]{\textcolor{White}{#1}}
\newcommand{\BT}[1]{\textcolor{Blue}{#1}}
\newcommand{\RT}[1]{\textcolor{red}{#1}}
\newcommand{\GT}[1]{\textcolor{green!45!yellow}{#1}}
%\newcommand{\TXT}{\rowcolor{lightgray}}



% Package microtype Warning: Unable to apply patch `footnote' 

\let\CheckCommand\providecommand
\usepackage{microtype}

% ---------------------------------
% Fancy Tables
% ---------------------------------

\usepackage{booktabs}
\usepackage{colortbl}
\usepackage{xfrac}
\newcommand{\AltLinha}[1]{\renewcommand{\arraystretch}{#1}}

\let\clipbox\relax
\usepackage{adjustbox}

% Animação LaTeX
\usepackage{animate}
\usepackage{graphics}

% Diagrama BNF
\usepackage{syntax}
%\setlength{\grammarparsep}{10pt plus 1pt minus 1pt} % increase separation between rules
%\setlength{\grammarindent}{10em} % increase separation between LHS/RHS 
\usepackage{tikz}

%\captionsetup{font=scriptsize}
\captionsetup{font={scriptsize,it}, % scriptsize + itálico, por exemplo
  %labelfont=bf,          % negrito na etiqueta (Figura 1, Tabela 2)
  labelsep=colon         % separador entre número e texto
}

\usepackage[font={scriptsize,it},skip=5pt]{caption}




\author{ Prof. Alberto Pena Lara}
\title{Aula No 08/15 : Plano de Trabalho}

\begin{document}

\section{Agenda}

\begin{frame}[plain]
\vspace{-2mm}
\maketitle
\vspace{-3mm}

\centering {\scriptsize{29 de abril de 2026}}\\

\vspace{4mm}

\begin{columns}

\begin{column}{0.3\textwidth}

\tcbset{left=2mm, top=2mm, right=2mm, bottom=2mm,center title}
\tcbset{toptitle=0.5mm, bottomtitle=0.5mm, boxrule=0.1mm, boxsep=0.15mm}
\tcbset{fonttitle=\scriptsize,colframe=blue!55!black}
\begin{tcolorbox}[fit to height=2cm,title= Agenda]
\begin{itemize} 
\itemsep1mm
\item Blocos logicos
\item O9erações
\item Estrutura Textual
\item Resumo
\end{itemize}
\end{tcolorbox}
\end{column}
\end{columns}
\end{frame}



\begin{frame}{}

\tikzstyle{naveqs} = [text width=70mm, fill=blue,  minimum height=20mm, rounded corners]
\tikzstyle{Aviso} = [text width=48mm, fill=green, minimum height=6mm, rounded corners, align=center] 
\tikzstyle{Visto} = [text width=48mm, fill=gray, text=white, minimum height=6mm, rounded corners, align=center] 


\begin{tikzpicture}
\node (naveq) [naveqs] {};
\path (naveq.center)+(0.2,0.55) node[anchor=center, xshift=8mm, text=white] {Elaboração do Projeto};
%\path (naveq.east)+(0.1, 0.5) node[anchor=west, text=blue] {\Large{\bf{Semiótica}}};
\path (naveq.east)+(0,-0.4) node (INS) [Aviso] {Pesquisa Semântica};
%\path (naveq.east)+(0, -1.1) node (INS) [greenRect,text width=68mm] {Estrutura Textual};
%\path (naveq.east)+(0, -1.8) node (INS) [greenRect,text width=48mm] {Exemplo concreto};
\end{tikzpicture}
\end{frame}



% ***********************************************

\begin{frame}{Conceitos}

\begin{columns}
\begin{column}{0.9\textwidth}

\begin{figure}[ht]
\centering
%\begin{tikzpicture}[>=Latex, line width=0.9pt,scale=0.5]
\begin{tikzpicture}[
    scale=0.5,
    transform shape,
    >=Latex,
    line width=0.9pt,
    every node/.style={font=\scriptsize}
]


\def\rx{0.20}
\def\ry{0.62}
\def\Lone{1.95}
\def\Ltwo{1.75}
\def\Lthr{1.65}
\def\Lfor{1.70}


% Cilindro 1
\coordinate (A1) at (0,0);
\draw (A1) ellipse ({\rx} and {\ry});
\draw (-\rx+0.2,\ry) -- ({\Lone-\rx},\ry);
\draw (-\rx+0.2,-\ry) -- ({\Lone-\rx},-\ry);
\draw ({\Lone-\rx},\ry) arc[start angle=90,end angle=-90,x radius={\rx},y radius={\ry}];

% Cilindro 2
\coordinate (A2) at (3.8,0);
\draw (A2) ellipse ({\rx} and {\ry});
\draw ({4-\rx},\ry) -- ({3.8+\Ltwo-\rx},\ry);
\draw ({4-\rx},-\ry) -- ({3.8+\Ltwo-\rx},-\ry);
\draw ({3.8+\Ltwo-\rx},\ry) arc[start angle=90,end angle=-90,x radius={\rx},y radius={\ry}];

% Cilindro 3
\coordinate (A3) at (7.2,0);
\draw (A3) ellipse ({\rx} and {\ry});
\draw ({7.4-\rx},\ry) -- ({7.2+\Lthr-\rx},\ry);
\draw ({7.4-\rx},-\ry) -- ({7.2+\Lthr-\rx},-\ry);
\draw ({7.2+\Lthr-\rx},\ry) arc[start angle=90,end angle=-90,x radius={\rx},y radius={\ry}];

% Cilindro 4
\coordinate (A4) at (10.5,0);
\draw[line width=1.2pt] (A4) ellipse ({\rx} and {\ry});
\draw[line width=1.2pt] ({10.7-\rx},\ry) -- ({10.5+\Lfor-\rx},\ry);
\draw[line width=1.2pt] ({10.7-\rx},-\ry) -- ({10.5+\Lfor-\rx},-\ry);
\draw[line width=1.2pt] ({10.5+\Lfor-\rx},\ry) arc[start angle=90,end angle=-90,x radius={\rx},y radius={\ry}];

% Conectores
\draw[->] ({\Lone+0.15},0) -- ({3.8-0.45},0);
\draw[->] ({3.8+\Ltwo+0.10},0) -- ({7.2-0.45},0);
\draw[->] ({7.2+\Lthr+0.10},0) -- ({10.5-0.45},0);

% Rótulos
\node[font=\Large] at (0.8,1.5) {Blocos lógicos};
\node[font=\Large, align=center] at (4.8,-1.8) {Pesquisa\\Semântica};
\node[font=\Large] at (7.9,1.5) {Artigos};
\node[font=\Large, align=center] at (11.4,-1.8) {Mineirar\\Dados};
\end{tikzpicture}
\vspace{5mm}
\caption{Representação esquemática do pipeline.}
\end{figure}

\end{column}
\end{columns}
\end{frame}



\begin{frame}{}

\tikzstyle{naveqs} = [text width=70mm, fill=blue,  minimum height=20mm, rounded corners]
\tikzstyle{Aviso} = [text width=48mm, fill=green, minimum height=6mm, rounded corners, align=center] 
\tikzstyle{Visto} = [text width=48mm, fill=gray, text=white, minimum height=6mm, rounded corners, align=center] 


\begin{tikzpicture}
\node (naveq) [naveqs] {};
\path (naveq.center)+(0.2,0.55) node[anchor=center, xshift=8mm, text=white] {Elaboração do Projeto};
%\path (naveq.east)+(0.1, 0.5) node[anchor=west, text=blue] {\Large{\bf{Semiótica}}};
\path (naveq.east)+(0,-0.4) node (INS)  [Visto] {Pesquisa Semântica};
\path (naveq.east)+(0, -1.1) node (INS) [Aviso] {Blocos Lógicos};
%\path (naveq.east)+(0, -1.8) node (INS) [greenRect,text width=48mm] {Exemplo concreto};
\end{tikzpicture}
\end{frame}


% ***********************************************

\begin{frame}{Blocos de Pesquisa}

\begin{columns}
\begin{column}{0.6\textwidth}

\tcbset{left=1mm, top=1mm, right=1mm, bottom=1mm}
\tcbset{toptitle=0.5mm, bottomtitle=0.5mm, boxrule=0.1mm, boxsep=0.15mm}
\tcbset{colframe=blue,coltitle=white}
\tcbset{fonttitle=\scriptsize,fontupper=\large}
\tcbset{segmentation style={solid, blue, opacity=1 ,line width=0.1mm}}

\begin{tcolorbox}[fit to height=3     cm, title= Ruído como perda de coerênc] 

(coherence OR decoherence OR "dephasing noise" OR "amplitude damping")
AND ("open quantum systems" OR Lindblad OR "quantum operations")

\tcbsubtitle{Relatividade como fonte de ruído efetivo}

("relativistic quantum information" OR "observer-dependent particles")
AND ("Unruh effect" OR "Bogoliubov transformations" OR "thermal state")
AND ("entanglement degradation" OR "non-inertial frames")

\end{tcolorbox}
\end{column}
\end{columns}
\end{frame}

% ***********************************************

\begin{frame}{Blocos de Pesquisa}

\begin{columns}
\begin{column}{0.6\textwidth}

\tcbset{left=1mm, top=1mm, right=1mm, bottom=1mm}
\tcbset{toptitle=0.5mm, bottomtitle=0.5mm, boxrule=0.1mm, boxsep=0.15mm}
\tcbset{colframe=blue,coltitle=white}
\tcbset{fonttitle=\scriptsize,fontupper=\large}
\tcbset{segmentation style={solid, blue, opacity=1 ,line width=0.1mm}}

\begin{tcolorbox}[fit to height=3cm, title= String composta recomendada] 

("relativistic quantum information" OR "quantum information in curved spacetime")
AND
("open quantum systems" OR "quantum operations" OR "Stinespring dilation" OR Lindblad)
AND
(decoherence OR noise OR thermalization OR "entanglement degradation")
AND
("Unruh effect" OR "Bogoliubov transformations" OR acceleration OR horizons OR Hawking)
\end{tcolorbox}
\end{column}
\end{columns}
\end{frame}

% ***********************************************

\begin{frame}{}

\tikzstyle{naveqs} = [text width=70mm, fill=blue,  minimum height=20mm, rounded corners]
\tikzstyle{Aviso} = [text width=48mm, fill=green, minimum height=6mm, rounded corners, align=center] 
\tikzstyle{Visto} = [text width=48mm, fill=gray, text=white, minimum height=6mm, rounded corners, align=center] 
\begin{tikzpicture}
\node (naveq) [naveqs] {};
\path (naveq.center)+(0.2,0.55) node[anchor=center, xshift=8mm, text=white] {Elaboração do Projeto};
%\path (naveq.east)+(0.1, 0.5) node[anchor=west, text=blue] {\Large{\bf{Semiótica}}};
\path (naveq.east)+(0,-0.4) node (INS)  [Visto] {Pesquisa Semântica};
\path (naveq.east)+(0, -1.1) node (INS) [Aviso] {Blocos Lógicos};
\path (naveq.east)+(0, -1.8) node (INS) [Aviso] {Operações};
\end{tikzpicture}
\end{frame}

\begin{frame}[fragile]{Resultados em CSV, Bibtex, Json}

\begin{columns}
\begin{column}{0.9\textwidth}

{\tiny
\begin{verbatim}
  @article{Blackburn2016,
  author    = {Blackburn, James A. and Cirillo, Matteo and Gr{\o}nbech-Jensen, Niels},
  title     = {A survey of classical and quantum interpretations of experiments on 
  Josephson junctions at very low temperatures},
  journal   = {Physics Reports},
  year      = {2016},
  volume    = {611},
  pages     = {1--34}
}

@article{Koch2007,
  author    = {Koch, Jens and Yu, Terri M. and Gambetta, Jay and Houck, A. A. and Schuster, D.
  I. and Majer, J. and Blais, Alexandre and Devoret, M. H. and Girvin, S. M. and Schoelkopf, R. J.},
  title     = {Charge-insensitive qubit design derived from the Cooper pair box},
  journal   = {Physical Review A},
  year      = {2007},
  volume    = {76},
  pages     = {042319}
}

@book{Nakahara2008,
  author    = {Nakahara, Mikio and Ohmi, Tetsuo},
  title     = {Quantum Computing: From Linear Algebra to Physical Realizations},
  year      = {2008},
  publisher = {CRC Press},
  address   = {Boca Raton},
  isbn      = {978-0-7503-0983-7}
}
\end{verbatim}
}
\end{column}
\end{columns}
\end{frame}


\begin{frame}{}

\tikzstyle{naveqs} = [text width=70mm, fill=blue,  minimum height=20mm, rounded corners]
\tikzstyle{Aviso} = [text width=48mm, fill=green, minimum height=6mm, rounded corners, align=center] 
\tikzstyle{Visto} = [text width=48mm, fill=gray, text=white, minimum height=6mm, rounded corners, align=center] 
\begin{tikzpicture}
\node (naveq) [naveqs] {};
\path (naveq.center)+(0.2,0.55) node[anchor=center, xshift=8mm, text=white] {Elaboração do Projeto};
%\path (naveq.east)+(0.1, 0.5) node[anchor=west, text=blue] {\Large{\bf{Semiótica}}};
\path (naveq.east)+(0,-0.4) node (INS)  [Visto] {Pesquisa Semântica};
\path (naveq.east)+(0, -1.1) node (INS) [Visto] {Blocos Lógicos};
\path (naveq.east)+(0, -1.8) node (INS) [Visto] {Operações};
\path (naveq.east)+(0, -2.5) node (INS) [Aviso] {Estrutura Textual};
\end{tikzpicture}
\end{frame}


% ***********************************************

\begin{frame}{Blocos de Pesquisa}

\begin{columns}
\begin{column}{0.3\textwidth}

\tcbset{left=1mm, top=1mm, right=1mm, bottom=1mm}
\tcbset{toptitle=0.5mm, bottomtitle=0.5mm, boxrule=0.1mm, boxsep=0.15mm}
\tcbset{colframe=blue,coltitle=white}
\tcbset{fonttitle=\scriptsize,fontupper=\large}
\tcbset{segmentation style={solid, blue, opacity=1 ,line width=0.1mm}}

\begin{tcolorbox}[fit to height=4cm, title= Artigo] 

\begin{itemize}
\itemsep1mm
\item Titulo
\item Subtítulo
\item Resumo
\item Palavras chaves
\item Introdução
\item Objetivos
\item Descrição do problema
\item Referencial Teórico
\item Resultados
\item Conclusões
\end{itemize}

\end{tcolorbox}
\end{column}
\end{columns}
\end{frame}





\begin{frame}
\frametitle{Metodologia}

\begin{columns}
\begin{column}{0.7\textwidth}

\tcbset{left=2mm, top=1mm, right=1mm, bottom=1mm,center title}
\tcbset{toptitle=0.5mm, bottomtitle=0.5mm, boxrule=0.1mm, boxsep=0.15mm}
\tcbset{fonttitle=\normalsize,colframe=blue,coltitle=white}
\tcbset{segmentation style={solid, blue, opacity=1 ,line width=0.1mm}}

\begin{tcolorbox}[fit to height=7cm,title={Resumo}]

\textcolor{red}{Contexto}. \textcolor{blue}{Gap}.  \textcolor{violet}{Propósito}.  \textcolor{green!80!black}{Resultados}. \textcolor{cyan}{Conclusões}

\tcbline


\textcolor{red}{A elaboração do Trabalho de Conclusão de Curso (TCC) representa um componente curricular obrigatório que visa consolidar competências investigativas e argumentativas dos estudantes no ensino superior. Contudo, a produção do TCC envolve múltiplas exigências metodológicas, cognitivas e operacionais, cujo atendimento satisfatório ainda se mostra um desafio significativo para a maioria dos discentes.} 
\textcolor{blue}{Apesar da ampla normatização institucional sobre a estrutura do TCC, observa-se carência de estudos que integrem, de forma sistemática, os diversos fatores objetivos que comprometem a qualidade dos trabalhos finais, especialmente no que se refere à articulação entre aspectos técnicos, argumentativos e operacionais.} \textcolor{violet}{Esta disciplina tem como objetivo identificar e descrever os principais fatores objetivos que dificultam a redação, apresentação e defesa do TCC, considerando tanto os aspectos estruturais da produção científica quanto as competências metodológicas exigidas ao longo do processo.} \textcolor{green!50!black}{A análise identificou nove categorias principais de dificuldades: delimitação temática inadequada; formulação imprecisa do problema de pesquisa; definição deficiente dos objetivos; uso de referencial teórico não científico; escolha metodológica incompatível; ausência de planejamento e cronograma; desconhecimento das normas acadêmicas; lacunas na coesão textual e argumentação científica; falhas na articulação entre teoria e dados.} \textcolor{cyan}{Tais dificuldades estão associadas a insuficiências formativas na escrita acadêmica, no raciocínio científico e na gestão do processo investigativo. Os dados indicam que o êxito na elaboração do TCC depende da integração de competências científicas, metodológicas e organizacionais, requerendo suporte institucional contínuo e estratégias pedagógicas específicas para o desenvolvimento da autonomia investigativa e da proficiência discursiva dos estudantes.}
\end{tcolorbox}
\end{column}
\end{columns}
\end{frame}


\begin{frame}[fragile]
\frametitle{}
\begin{Constative}{A-06}{Avaliação}{blue}
Aquecimento no desenvolvimento do TCC
\vspace{5mm}
\begin{enumerate}[label=(\alph*),start=1]
\itemsep1mm
\item Quais são os principais elementos textuais de um artigo científico, segundo a aula?
% Resposta:** A aula organiza a estrutura textual do artigo em: **título, subtítulo, resumo, palavras-chave, introdução, objetivos, descrição do problema, referencial teórico, resultados e conclusões**. Essa sequência mostra que a escrita científica não é apenas redação livre, mas uma composição formal orientada por funções discursivas bem definidas. 
\item Como o resumo deve ser estruturado metodologicamente?
% Resposta: O resumo deve seguir a sequência: **contexto, gap, propósito, resultados e conclusões**. Em termos científicos, isso significa: primeiro situar o tema; depois explicitar a lacuna de conhecimento; em seguida declarar o objetivo do estudo; posteriormente sintetizar os principais achados; e, por fim, enunciar a conclusão central. Essa organização torna o resumo semanticamente coerente e metodologicamente legível. 

\item O que significa organizar uma pesquisa por “blocos lógicos”?
% Resposta:Significa decompor o problema de pesquisa em **núcleos conceituais articulados**, convertendo cada núcleo em descritores de busca. Na aula, isso aparece como estratégia de **pesquisa semântica**, em que conceitos relacionados são agrupados para formar consultas mais robustas. Em vez de buscar palavras soltas, o pesquisador estrutura o problema em blocos conceituais, o que aumenta precisão, rastreabilidade e consistência metodológica. 

\item Como se constrói uma string booleana coerente para pesquisa semântica?
% Resposta: A construção segue uma lógica combinatória: * usa-se **OR** para reunir sinônimos, variantes ou termos próximos dentro do mesmo bloco conceitual;    * usa-se **AND** para conectar blocos conceituais diferentes que precisam aparecer simultaneamente; eventualmente, usam-se expressões entre aspas para preservar termos compostos. A aula exemplifica isso ao combinar blocos como **informação quântica relativística**, **sistemas quânticos abertos**, **ruído/decoerência** e **efeito Unruh/Bogoliubov/Hawking** em uma única string composta. 

\item Quais questões a aula propõe para avaliar o aquecimento do desenvolvimento do TCC?
% Resposta:** A aula propõe questões operacionais centrais, tais como:    * como transformar uma ideia ampla em **tema, objeto, problema e dados** bem delimitados;  * como converter uma pergunta em linguagem comum para **linguagem científica adequada à pesquisa semântica**;  * como mapear descritores, organizar blocos lógicos e construir uma **string booleana coerente**; quais critérios indicam que a pesquisa semântica está bem construída;  * como integrar **metodologia, evidências, ética e LGPD** quando a pesquisa envolve dados de pessoas.      Essas perguntas funcionam como matriz diagnóstica do amadurecimento metodológico do projeto. 

\end{enumerate}
\end{Constative}
\end{frame}

\end{document}
